\documentclass[12pt]{article}

% Idioma y codificación
\usepackage[spanish]{babel}
\usepackage[utf8]{inputenc}

% Márgenes
\usepackage{geometry}
\geometry{margin=2.5cm}

% Paquetes útiles
\usepackage{graphicx} % Imágenes
\usepackage{listings} % Código
\usepackage{xcolor}

% Configuración para código en C
\lstset{
    language=C,
    basicstyle=\ttfamily\small,
    numbers=left,
    numberstyle=\tiny,
    frame=single,
    breaklines=true
}

\begin{document}

% ---------- PORTADA ----------
\begin{titlepage}
    \centering
    \vspace*{4cm}
    
    {\Large \textbf{Programación Estructurada}}\\[1cm]
    {\large Práctica No. 2}\\[2cm]
    
    {\large Alexis Roberto Constantino Núñez}\\
		
    {\large Matricula: 2201619}\\
    {\large Grupo 531}\\[1cm]
    
    {\large \today}
\end{titlepage}

% ---------- PROBLEMA 1 ----------
\section{Problema 1}

\subsection{Enunciado del problema}
%Aquí va una breve explicación del programa realizado en lenguaje C.  
%Máximo 5 líneas explicando qué hace el programa y cómo funciona.
Suponga que existe un terreno cuyo ancho es tres veces mayor que el largo. Se pide
diseñe e implemente un algoritmo que lea solo el ancho del terreno y con ese dato
calcule el perímetro y el área del terreno.
\subsection{Especificación de las entradas y salidas}
Entradas: Ancho del terreno.\\
Salidas: Perímetro y Área del terreno.
\subsection{Algoritmo del programa}
\begin{verbatim}
PREGUNTAR ancho
CALCULAR LARGO l = ancho / 3
CALCULAR ÁREA a = ancho * largo
CALCULAR PERÍMETRO p = 2 * (ancho + largo)
IMPRIMIR área y perímetro
\end{verbatim}
\subsection{Pruebas de escritorio }
\begin{verbatim}
  MENU
Opcion 1: Calcular el area y perimetro de un terreno.
Opcion 2: Area de un triangulo.
Opcion 3: Salir.
 1
Calculo de un terreno cuyo ancho es tres veces el largo.
 Escriba el ancho del terreno:
57.83

El area del terreno es de 1114.7697 m^2.                                                       
\end{verbatim}

\subsection{Implementacion del algoritmo en C}
\begin{lstlisting}
/*Problema 1*/
void terreno (void)
{
        float ancho, largo;
        printf("Calculo de un terreno cuyo ancho es tres veces el largo.\n ");
        printf("Escriba el ancho del terreno:\n");
        scanf(" %f", &ancho);
        largo = ancho / 3;
        printf("\nEl area del terreno es de %.4f m^2.\n", farea(ancho, largo));
        printf("\nEl perimetro del terreno es de %.4f m.\n", fperimetro(ancho, largo));
}
//area del terreno
float farea (float ancho, float largo)
{
        float area;
        area = ancho * largo;
        return area;
}
//perimetro del terreno
float fperimetro (float ancho, float largo)
{
        float perimetro;
        perimetro = 2 * (ancho + largo);
        return perimetro;
}
\end{lstlisting}
% ---------- PROBLEMA 2 ----------
\section{Problema 2}

\subsection{Enunciado del problema}
%Aquí va una breve explicación del programa realizado en lenguaje C.  
%Máximo 5 líneas explicando qué hace el programa y cómo funciona.
Diseñe e implemente un algoritmo para calcular el área de un triángulo mediante
la fórmula:
\[
	A = \sqrt{p(p-a)(p-b)(p-c)}
\]
donde p es el semiperimetro: 
\[
	p = \frac{a + b + c}{2}
\]
\subsection{Especificación de las entradas y salidas}
Entradas: Primer, Segundo y Tercer lado del triangulo.\\
Salidas: Área del triangulo.
\subsection{Algoritmo del programa}
\begin{verbatim}
PREGUNTAR lado1
PREGUNTAR lado2
PREGUNTAR lad3o
CALCULAR SEMIPERIMETRO p = (a + b + c) / 2
CALCULAR ÁREA a = [{p * (p-a) * (p - b) * (p - c)}] ^ 1/2
IMPRIMIR área
\end{verbatim}
\subsection{Pruebas de escritorio }
\begin{verbatim}
   MENU
Opcion 1: Calcular el area y perimetro de un terreno.
Opcion 2: Area de un triangulo.
Opcion 3: Salir.
 2
Calculo del area de un triangulo por la formula de Heron.
Lado A:
47.2
Lado B:
88
Lado C:
126.9
El area del triangulo es: 131.0500
                MENU
Opcion 1: Calcular el area y perimetro de un terreno.
Opcion 2: Area de un triangulo.
Opcion 3: Salir.
 3
Saliendo...
\end{verbatim}

\subsection{Implementacion del algoritmo en C}
\begin{lstlisting}
/*Problema 2*/
void heron (void)
{
        float a, b, c;
        float smper, area;
        printf("Calculo del area de un triangulo por la formula de Heron.\n");
        printf("Lado A:\n"); scanf(" %f", &a);
        printf("Lado B:\n"); scanf(" %f", &b);
        printf("Lado C:\n"); scanf(" %f", &c);
        smper = semiperimetro (a, b, c);
        area = areatrang (smper, a, b, c);
        printf("El area del triangulo es: %.4f\n", area);
}
//perimetro del triangulo
float semiperimetro (float a, float b, float c)
{
        float sperimetro;
        sperimetro = (a+b+c)/2;
        return sperimetro;
}
//area del triangulo
float areatrang (float p, float a, float b, float c)
{
        float area, aux;
        aux = p * (p-a)*(p-b)*(p-c);
        area = pow (aux, 0.5);
        return area;
}

\end{lstlisting}


\end{document}

